\begin{center}
\begin{minipage}{27pc}
  \small
  \noindent{\scshape abstract.}\enspace
  It is often suggested that no decision theory can be optimal on every fair decision problem,
  because the standards by which theories are judged conflict. This paper makes the claim precise
  for three successively stronger notions of a decision theory and determines when it holds. If a
  theory is any assignment of policies to finite problems whose structural representation supplies
  every counterfactual, ex-ante optimality is always attainable, and every inconsistency among
  ex-ante optimality, causal dominance, ratifiability, commitment stability, and an evidential
  criterion has a core of one or two criteria; the conflict between ex-ante optimality and dominance
  is possible exactly when the state law depends on the prescribed action. If a theory is a rule on
  causal matrices, whose rows are the belief states induced by the agent's policies and whose
  diagonal is ex-ante value, common affine invariance again yields only small cores, while
  invariance under row translations, with neutrality and contraction, forces a rule to choose every
  option on exactly the domains satisfying a reachability condition, and principal binary
  independence characterises ex-ante evaluation; both results fail if a smaller menu keeps the belief
  model of the larger one. If a theory is a computable procedure that reads the source code of its
  environment and must infer the dependence itself, a uniform optimiser exists under a designated
  substitution semantics and a termination promise; the negative results come from fixed copies of
  the agent, from the uncomputability of exact selection, and from the ways of recognising the agent
  in code. An axiomatisation of counterfactual operators shows that calibration together with an
  extensionality that identifies computations by their actual behaviour forces causal evaluation and
  is inconsistent with ratifiability, while an extensionality that compares computations under every
  substitution admits ex-ante evaluation satisfying every desideratum. Each impossibility concerns how
  the agent's alternatives are specified, not how to choose among them.
\end{minipage}
\end{center}
